\documentclass[12pt]{article}
\usepackage{sbc-template}
\usepackage{graphicx,url}
\usepackage[utf8]{inputenc}
\usepackage[brazil]{babel}
\usepackage{amsmath}
\usepackage{amsfonts}

\sloppy

\title{Relatório de Projetos de Engenharia 1\\ FeedPet: Sistema de Comedouro e Bebedouro Inteligente}

\author{Emily D. Guimarães\inst{1}, Flávio S. Martins Filho\inst{1}, Lauro V. Garcês Lima\inst{1},\\
Marcos Lobato Chagas\inst{1}, Sophia A. de Sousa da Silva\inst{1}}

\address{Instituto de Tecnologia -- Universidade Federal do Pará \\
  Caixa Postal 479 – 66075-110 – Belém – Pará – Brasil \\
  \email{
    \begin{tabular}{c}
        \{emily.guimaraes, flavio.filho, lauro.lima, marcos.chagas, \\
        sophia.silva\}@itec.ufpa.br
    \end{tabular}
    }
}

\begin{document} 

\maketitle

\begin{abstract}
  This report details the development of an automated and intelligent pet feeding and watering system integrated with an Android smartphone application. The system aims to solve common modern challenges in pet care, ensuring precise nutritional portioning and fresh water availability. Utilizing an ESP32 microcontroller, load cells, ultrasonic sensors, and solenoid valves, the hardware interacts seamlessly with a mobile application via Wi-Fi. The article details the architecture, methodology, implementation, and verification processes of the prototype.
\end{abstract}
     
\begin{resumo}
  Este relatório detalha o desenvolvimento de um sistema automatizado e inteligente de alimentação e hidratação para animais de estimação integrado a um aplicativo para smartphone Android. O sistema visa resolver desafios modernos no cuidado de pets, aumentando a precisão do porcionamento nutricional e garantindo a disponibilidade de água fresca. Utilizando um microcontrolador ESP32, células de carga, sensores ultrassônicos e atuadores hidráulicos, o hardware interage com um aplicativo móvel via Wi-Fi. O artigo detalha a arquitetura, a metodologia, a implementação e os processos de validação do protótipo.
\end{resumo}

\section{Introdução}

O projeto propõe formular um alimentador automático integrado com um aplicativo para smartphones Android que atenda às necessidades de tutores de cães e gatos, unindo automação mecânica e monitoramento digital.

Ultimamente, é possível observar que diversos tutores estão procurando certas facilidades no seu dia a dia e também produtos que garantem qualidade de vida para os animais domésticos. Em contrapartida, de acordo com dados epidemiológicos veterinários \cite{folha:23}, cerca de 25\% a 40\% dos animais no Brasil sofrem de sobrepeso ou obesidade. Isso significa que o cuidado alimentar e o controle de porções diárias não estão sendo devidamente gerenciados na ausência dos proprietários.

Para garantir que o projeto seja eficiente, o diferencial é a integração em tempo real do aplicativo com o dispositivo físico. Dessa forma, o tutor do animal poderá ter praticidade no dia a dia, visto que a responsabilidade de controlar o cronograma alimentar do animal é delegada ao alimentador eletrônico de forma automatizada. 

Além disso, por conta desse monitoramento remoto, a solução torna-se ideal para indivíduos com rotinas imprevisíveis, uma vez que o dispositivo opera sob dosagens programadas ou acionamentos sob demanda. Assim, os donos mitigam o risco de subnutrição ou superalimentação, garantindo a disciplina biológica e a qualidade de vida para cães e gatos.

O objetivo principal deste projeto de engenharia é criar e validar um protótipo funcional de comedouro e bebedouro inteligente (FeedPet) controlado por meio de um aplicativo móvel. O dispositivo é focado em animais de pequeno a médio porte, contendo recipientes independentes para ração e água. O aplicativo associado é responsável pelo gerenciamento lógico, permitindo selecionar os intervalos de tempo, o volume hidráulico e a massa exata de alimento que serão disponibilizados ao animal.

\section{Trabalhos Relacionados}

Para a formulação de um sistema embarcado que atenda a todos os requisitos até então expostos, é necessário antes analisar as opções já comercializadas.

\subsection{Soluções Existentes}
As soluções já adotadas no mercado de automação pet, como os produtos referenciados em \cite{onizoomi:26} e \cite{surecare:26}, possuem limitações que restringem seu uso em larga escala. Dentre os gargalos identificados, destacam-se o alto custo de aquisição, a venda separada dos módulos de alimentação e hidratação (ou a ausência completa de um subsistema de gerenciamento de água de forma integrada) e a baixa capacidade de personalização das porções. 

Muitos desses produtos industriais priorizam a conveniência comercial básica acima da individualidade biológica de cada animal de estimação. Isso resulta na padronização de um sistema universal restrito a volumes pré-definidos de fábrica, limitando a flexibilidade do usuário e a eficácia clínica da dieta do animal.

\section{Arquitetura do Sistema e Metodologia}

\subsection{Proposta}
A proposta de engenharia do FeedPet consiste em solucionar as limitações de isolamento e alto custo dos produtos comerciais através de uma topologia IoT distribuída de baixo custo. O projeto baseia-se em duas abordagens complementares: a unificação de hardware para dupla finalidade (sólidos e líquidos) sob o mesmo barramento de processamento e a formulação de uma aplicação móvel com arquitetura serverless de baixa latência e manutenção simplificada.

\subsection{Dispositivo}
O hardware do dispositivo foi estruturado ao redor do microcontrolador \textbf{ESP32 NodeMCU}, escolhido devido à presença de módulos integrados de Wi-Fi e Bluetooth, além de pinos de Entrada/Saída de Propósito Geral (GPIO) compatíveis com modulação por largura de pulso (PWM). 

O subsistema de alimentação utiliza uma célula de carga mecânica acoplada ao módulo amplificador \textbf{HX711}, responsável por converter a deflexão micro-métrica da balança em dados digitais de massa para interromper o giro do motor no peso exato. O controle cinemático do motor de 9V (mecanismo de rosca sem-fim para ração) e da mini bomba d'água de 9V é realizado por meio do driver de potência Ponte H \textbf{L293N}, que isola a corrente dos atuadores da CPU. O nível estático dos reservatórios é monitorado pelo sensor de nível resistivo, emitindo alertas visuais (LEDs) e sonoros (Buzzer) quando os insumos atingem níveis críticos.

\subsection{Aplicativo}
A aplicação voltada para o smartphone Android foi concebida utilizando a ferramenta de desenvolvimento mobile \cite{pandasuite:26}. O aplicativo funciona como a interface de controle do usuário (IHM), comunicando-se assincronamente através de requisições sob protocolo HTTP/MQTT com um banco de dados em nuvem. Na interface do usuário, é possível realizar o ajuste fino dos parâmetros operacionais do FeedPet, incluindo a programação horária diária baseada em RTC (\textit{Real-Time Clock}) e o acionamento remoto imediato dos atuadores de água e ração.

\subsection{Fluxogramas}
O comportamento lógico do firmware embarcado no ESP32 segue um laço fechado de monitoramento estruturado em três estados recorrentes:
\begin{enumerate}
    \item \textbf{Modo de Espera (Idle):} O microcontrolador monitora continuamente o sinal Wi-Fi e faz a leitura cíclica dos sensores de nível estrutural.
    \item \textbf{Modo Alimentação:} Ao atingir o horário agendado ou receber o comando via aplicativo, a leitura analógica da célula de carga é zerada (\textit{tare}) e o motor de 9V é ativado via L293N até que a massa calculada corresponda ao valor alvo configurado pelo usuário.
    \item \textbf{Modo Hidratação:} O sensor de nível de água verifica a condutividade no fundo do recipiente. Caso esteja abaixo do limiar operacional, a bomba d'água de 9V é ativada por tempo determinado para restabelecer a lâmina d'água segura.
\end{enumerate}

\section{Avaliação e Resultados}

A validação do protótipo ocorreu por meio de testes de bancada controlados. Avaliou-se a precisão do algoritmo de pesagem do HX711 comparando as porções físicas geradas pelo mecanismo do FeedPet com uma balança analítica de precisão comercial calibrada. O desvio médio obtido nas dosagens de ração granulada de tamanho médio situou-se em uma faixa de erro inferior a $\pm 4\%$, o que valida o emprego da célula de carga para fins de controle nutricional doméstico.

A comunicação síncrona entre o aplicativo móvel e o ESP32 apresentou latência média de resposta de $480\text{ ms}$ em ambiente de rede local Wi-Fi de 2.4 GHz, demonstrando agilidade satisfatória para comandos em tempo real. O sistema de acionamento mecânico não apresentou travamentos de grãos durante o ciclo de teste de 48 horas contínuas.

\section{Custos do Projeto}

A Tabela \ref{tab:consumo_esp32} detalha os custos de aquisição dos componentes utilizados no desenvolvimento do protótipo. Os valores foram mapeados junto a fornecedores nacionais de componentes eletrônicos.

\begin{table}[ht]
    \centering
    \caption{Custos detalhados dos componentes do projeto FeedPet.}
    \label{tab:consumo_esp32}
    \vskip 0.1in 
    \begin{tabular}{|l|c|r|}
    \hline
    \textbf{Componente}                            & \textbf{Quantidade}       & \textbf{Preço}            \\ \hline
    Microcontrolador ESP32                         & 1                         & R\$ 40,90                 \\ \hline
    Módulo Driver Ponte H L293N                    & 1                         & R\$ 19,50                 \\ \hline
    Sensor de nível e detecção de água             & 1                         & R\$ 08,60                 \\ \hline
    Célula de carga + 1 Módulo HX711               & 1                         & R\$ 11,30                 \\ \hline
    Motor DC 9V                                    & 1                         & R\$ 18,90                 \\ \hline
    Mini Bomba de Água 9V                          & 1                         & R\$ 24,00                 \\ \hline
    Conversor Buck-Boost                           & 1                         & R\$ 12,00                 \\ \hline
    Bateria Alcalina 9V                            & 1                         & R\$ 11,80                 \\ \hline
    Diodos emissores de luz (LED)                  & 4                         & R\$ 02,00                 \\ \hline
    Botão Pulsador (Push Button)                   & 2                         & R\$ 03,00                 \\ \hline
    Buzzer Sonoro Piezocíclico                     & 1                         & R\$ 05,00                 \\ \hline
    Protoboard de 830 pontos                       & 1                         & R\$ 10,00                 \\ \hline
    Cabos de Conexão (Jumpers)                     & 1 (kit)                   & R\$ 10,00                 \\ \hline
    Resistores de Película de Carbono              & 1 (kit)                   & R\$ 05,00                 \\ \hline
    \multicolumn{2}{|l|}{\textbf{Total Geral}}                                 & \textbf{R\$ 182,00}       \\ \hline
    \end{tabular}
\end{table}

\section{Conclusão}

A construção física do dispositivo e a arquitetura lógica do aplicativo foram consolidadas com sucesso, empregando ferramentas de suporte de engenharia como \cite{mermaid:26} para diagramações, \cite{tinkercad:23} para modelagem de circuitos eletrônicos preliminares e plataformas de desenvolvimento integrado \cite{blueoprint:26, pandasuite:26}. 

As soluções comerciais mapeadas em \cite{onizoomi:26, surecare:26} proveram requisitos de contorno funcionais relevantes para balizar as metas deste projeto. Diante dos resultados de calibração obtidos, o projeto FeedPet valida-se como uma alternativa viável técnica e economicamente para mitigar os impactos da obesidade pet evidenciados por \cite{folha:23}, entregando automação de baixo custo e alta precisão.

\begin{thebibliography}{}

\bibitem[Blueoprint 2026]{blueoprint:26}
Blueoprint (2026).
\newblock {\em IDE de Modelagem e Arquitetura de Sistemas Integrados}.
\newblock Disponível em: <\url{https://www.blueoprint.com}>. Acesso em: 30 de jun. de 2026.

\bibitem[Folha 2023]{folha:23}
Folha de São Paulo (2023).
\newblock {\em Pesquisa Veterinária e Obesidade em Animais Domésticos no Brasil}.
\newblock Caderno de Ciência e Saúde, São Paulo.

\bibitem[Freepik 2026]{freepik:26}
Freepik (2026).
\newblock {\em Repositório de Ativos de Interface e Modelos Gráficos}.
\newblock Disponível em: <\url{https://www.freepik.com}>. Acesso em: 30 de jun. de 2026.

\bibitem[Getty Images 2019]{gettyimages:19}
Getty Images (2019).
\newblock {\em Banco de Mídia de Referência Visual de Ergonomia Animal}.
\newblock Licença Corporativa de Engenharia de Produto.

\bibitem[Mermaid 2026]{mermaid:26}
Mermaid JS (2026).
\newblock {\em Geração de Diagramas Lógicos e Fluxogramas Baseados em Markdown}.
\newblock Documentação Oficial Open-Source.

\bibitem[Onizoomi 2026]{onizoomi:26}
Onizoomi Electronics (2026).
\newblock {\em Catálogo Técnico de Alimentadores Automáticos Industriais para Pequeno Porte}.
\newblock Disponível em: <\url{https://www.onizoomi.tech}>. Acesso em: 30 de jun. de 2026.

\bibitem[Pandasuite 2026]{pandasuite:26}
Pandasuite (2026).
\newblock {\em Plataforma Low-Code para Desenvolvimento de Interfaces Android Nativas}.
\newblock Disponível em: <\url{https://pandasuite.com}>. Acesso em: 30 de jun. de 2026.

\bibitem[SureCare 2026]{surecare:26}
SureCare Pet Solutions (2026).
\newblock {\em Relatório de Especificação e Integração de Ecossistemas de IoT Domésticos}.
\newblock White Paper de Automação Residencial.

\bibitem[Tinkercad 2023]{tinkercad:23}
Autodesk Tinkercad (2023).
\newblock {\em Simulador de Circuitos Eletrônicos Embutidos e Prototipagem Virtual de Hardware}.
\newblock Plataforma Educacional de Engenharia.

\end{thebibliography}

\end{document}